\section{评价量与四格比较}
\label{sec:framework}

\subsection{候选、标准、环境与分数}
\label{subsec:score}

考虑这样一个程序任务：输入两个布尔值，当且仅当两者不同时输出真。一个候选程序总是输出假，另一个正确计算异或。两者都可能通过 Lean 编译，但只有后者在每种输入下都给出所需答案。因此，给出分数之前，需要先说明它衡量的是哪一种性质。

用 $C$ 表示候选成果，包括保存的代码和事先选定的检查目标。这个目标称为声明，也就是代码中有名字的定义或定理；预先选定它，便固定了评价对象。用 $R$ 表示审核标准，规定做什么检查以及何时接受；用 $E$ 表示执行环境，包括编译器、库、搜索路径和运行设置。\emph{评价程序} $g$ 则在环境 $E$ 中，按标准 $R$ 检查候选 $C$。

当一次检查确实取得了有效的数值结果时，我们把这个分数写成
\[
q=g(C,R;E).
\]
这里的 $q$ 就是指定程序返回的分数。比较保持环境不变。本报告使用的接受分数没有量纲：$1$ 表示指定检查通过，$0$ 表示未通过。它们是这项检查的判断，并不是数学内容正确的概率。

使用 $g$ 时，仍需说明实际程序、检查目标、标准和证据；一个函数符号本身并不能定义通用的数学正确性评分。下面先看本次核查实际实现的程序。

\subsection{公理政策测量及缺失、无效分数}
\label{subsec:measurement-status}

发布包中的 \code{evaluate\_sources.py} 实现的是范围较窄的公理政策检查。
它保存源代码，为事先选定的声明追加查询命令，调用已记录的 Lean 4.31.0 编译器，
然后提取两项信息：声明打印出来的类型，以及系统报告的公理依赖。
类型描述这个声明的接口或命题；公理列表记录这个声明依赖的假设。
这些观察很有用，但任何一项都不会自动说明原任务究竟想表达什么。

在这个程序中，标准 $R$ 提供允许使用的公理名单。当且仅当提取出的每个公理都在允许名单中时，
分数为 $1$；只要发现至少一个不允许的公理，分数就是 $0$。
例如，若已成功提取的依赖中出现 \code{sorryAx}，而标准禁止它，则得到 $0$。
空的公理依赖列表可以通过这样的标准，但“没有获得编译器输出”不等于“依赖列表为空”。
Lean 的文档同样区分证明的有效性与被证明命题的含义。\cite{leanvalidation}

在比较分数之前，程序还会用固定打印设置检查两份声明的类型文本是否完全相同。这是一个保守的接口匹配条件。它仍不能确定行为：前面的两个布尔程序都可以具有 \code{Bool -> Bool -> Bool} 类型。打印类型不同时，本实现不建立数值比较，但这并不排除另行设计有意义的比较。

若要检查行为，可以把四种输入全部运行一遍，再与题目要求的真值表比较。后面的有限域对照正是增加了这项测试，从而取得单靠公理政策检查无法提供的证据。

检查也可能没有产生可用分数，因此还需要保留测量状态：

\begin{description}[leftmargin=1.2em,style=nextline]
\item[已测量。] 指定检查确实得到了数值。已测量的零分表示这项检查实际拒绝了该候选成果。
\item[未运行、材料缺失或结果不确定。] 想问的问题有意义，但没有取得答案。
例如，编译器不可用、源文件缺失、运行超时，或者没有成功提取选中的声明。
这个范围较窄的公理检查不会把编译失败自动记为公理政策零分，因为它并没有取得公理测量结果。
\item[无效或不适用。] 所提议的比较缺少必要的共同含义，或者这项标准不适用于选定对象。
若把这种格子当成一个尚未知晓的数值，就已经偷偷假定了一个尚未被证明合理的比较。
\end{description}

怎样编码取决于所问的问题。研究编译成功与否时，可以把编译失败记为零；公理政策检查尚未取得公理测量时，则不能作同样处理。

\subsection{四格、条件差与交互量}
\label{subsec:four-cells}

现在固定评价程序和执行环境，在两个标准下分别比较两个候选成果。
其中，$C_0$ 和 $C_1$ 是两个候选成果，$R_0$ 和 $R_1$ 是两个审核标准。记
\[
q_{ij}=g(C_i,R_j;E),\qquad i,j\in\{0,1\}.
\]
下面的例子给出四种候选与标准组合的分数。

下面是一个\emph{用于说明计算的例子，不是新做的实验}。
仍然使用刚才的两个布尔程序：$C_0$ 总是输出假，$C_1$ 则在两个输入不同时输出真。
让 $R_0$ 只检查输入“假、假”，该输入的要求输出是假；让 $R_1$ 检查全部四种输入。
按“假、假”“假、真”“真、假”“真、真”的顺序，要求的输出依次是假、真、真、假。
这个例子里的 $g$ 表示：“运行标准指定的输入；所有输出都符合要求就给 $1$，否则给 $0$。”
它是为这个例子定义的行为测试程序，不是发布包中只测公理政策的程序。

\begin{center}
\begin{tabular}{lcc}
\toprule
候选成果 & $R_0$：一个输入 & $R_1$：全部四个输入\\
\midrule
$C_0$：总是输出假 & $q_{00}=1$ & $q_{01}=0$\\
$C_1$：计算异或 & $q_{10}=1$ & $q_{11}=1$\\
\bottomrule
\end{tabular}
\end{center}

先保持审核标准不变，比较两个候选成果。用 $D_{C\mid R_0}$ 表示在 $R_0$ 下，
第二个候选成果的分数减去第一个候选成果的分数。
$D_{C\mid R_1}$ 则是在 $R_1$ 下作同样的比较。沿着两列表格分别向下读，得到
\begin{align*}
D_{C\mid R_0}&=q_{10}-q_{00}=1-1=0,\\
D_{C\mid R_1}&=q_{11}-q_{01}=1-0=1.
\end{align*}
只测一个输入时，检查分不出这两个程序；测遍四种输入时，就能分出来。
所以，报告中的候选成果差异，取决于比较时固定的是哪一个标准。

接下来保持候选成果不变，比较两个标准。沿着每一行分别向右读，得到
\begin{align*}
D_{R\mid C_0}&=q_{01}-q_{00}=0-1=-1,\\
D_{R\mid C_1}&=q_{11}-q_{10}=1-1=0.
\end{align*}
扩大后的测试拒绝了错误程序，但正确程序仍然通过。
第一个差值为负，只表示测试得分下降，不能解释成“标准改进使程序本身变差了”。
这四个差值的单位都是\emph{分数点}，不是“程序质量的百分比”。

如果沿对角线比较，就同时改变了两项要素。用 $\Delta$ 表示这个总变化：
\[
\Delta=q_{11}-q_{00}=1-1=0.
\]
如果只看这两个接受标签，就会同时漏掉程序的改进与测试要求的提高。
\emph{交互量}记作 $I$，表示标准改变后，候选成果差异改变了多少：
\begin{align*}
I&=D_{C\mid R_1}-D_{C\mid R_0}=1-0=1\\
 &=q_{11}-q_{10}-q_{01}+q_{00}.
\end{align*}
把同样四项重新分组，还可得到
$I=D_{R\mid C_1}-D_{R\mid C_0}$。
这里的交互量是表格的算术描述，并不证明存在某种心理上的合作或因果机制。
上述四个条件差、一个对角差和一个交互量，就是本次审计所使用的六个诊断问题。

\subsection{两要素夏普利分摊}
\label{subsec:allocation}

从左上格走到右下格有两条路径：可以先换候选成果，再换标准；也可以先换标准，再换候选成果。
无论走哪条路径，两次变化的和都等于 $\Delta$。
如果我们决定对这两种顺序给予相同权重，就把候选成果的两个条件差取平均，作为分给候选成果的变化量；
再把标准的两个条件差取平均，作为分给标准的变化量。
两个分摊量分别记作 $\phi_C$ 和 $\phi_R$：
\begin{align*}
\phi_C&=\tfrac12[(q_{10}-q_{00})+(q_{11}-q_{01})],\\
\phi_R&=\tfrac12[(q_{01}-q_{00})+(q_{11}-q_{10})].
\end{align*}
这就是标准的两要素夏普利分摊。\cite{shapley}
在上例中，$\phi_C=(0+1)/2=1/2$，$\phi_R=(-1+0)/2=-1/2$。
它们的和为零，与对角线总变化一致。
一般情况下，把两个式子相加，中间各项抵消，得到
$\phi_C+\phi_R=q_{11}-q_{00}=\Delta$。

这种分摊以分数点为单位，取决于两种顺序等权的约定。它概括的是表中已有的比较，并非从历史中恢复了因果贡献。交互量已经包含在平均数内，再加一次就会破坏总和等式。

\subsection{历史审核：记录标签与部分可观察的条件}
\label{subsec:historical-model}

四格比较中的程序和标准是明确规定的。历史审核则较难控制：档案可能保存了判断，却只保留了产生判断的部分过程。对一次审核事件 $e$，需要区分以下角色：
\begin{itemize}
\item $Y_e$ 是记录中的结果标签，例如通过或不通过。在另行说明数值编码之前，它只是类别。
\item $O_e$ 是当时实际接受审查的对象，包括候选代码以及相关上下文。
\item $K_e$ 是那次审查中实际采用的有效标准。
\item $M_e$ 是审查装置或设置，例如审查者或模型、提示文本、工具以及运行配置。
\item $U_e$ 汇集这套描述没有分别记录的其余影响。这里并没有为它假定概率分布，也没有假定它与其他量独立。
\end{itemize}
于是，概念式
\[
Y_e=\psi(O_e,K_e,M_e,U_e)
\]
只是说：这些要素通过某种审查过程产生了结果；这个过程用 $\psi$ 表示。
它不像已经拟合的回归方程那样，提供一组可以估计的系数，也没有给出已知的 $Y_e$ 预测规则。
早期记号用 $C_e^*$ 表示有效标准，用 $\varepsilon_e$ 表示其余影响；
这里改用 $K_e$ 与 $U_e$，是为了避免把那里的标准与四格表中的候选成果 $C$ 混淆。

例如，合同文本的哈希值可以标识保存了哪份文字，却不能确定审核者在判断时怎样执行了它。实际标准仍可能只被部分观察到。同样，标签变化可能伴随候选、要求或审核程序的变化。后面的过程分析按明确路径规则研究记录序列，并不能单独识别这些影响。

不过，我们仍能研究：一项明确的测量能够区分什么，一个给定的观察模型会推出什么，以及现有证据能确定哪些比较。